\documentclass[11pt,oneside]{article}

\usepackage[a4paper,margin=27mm,headheight=15pt]{geometry}
\usepackage[T1]{fontenc}
\usepackage[utf8]{inputenc}
\IfFileExists{libertinus.sty}{\usepackage{libertinus}}{\usepackage{lmodern}}
\usepackage{microtype}
\usepackage{amsmath,amssymb,amsthm,mathtools,mathrsfs}
\usepackage{bm}
\usepackage{booktabs,longtable,array,tabularx}
\usepackage{graphicx}
\usepackage{xcolor}
\usepackage{enumitem}
\usepackage{fancyhdr}
\usepackage{titlesec}
\usepackage{listings}
\usepackage{xurl}
\usepackage{hyperref}
\usepackage[nameinlink,noabbrev]{cleveref}

\allowdisplaybreaks
\numberwithin{equation}{section}
\setcounter{tocdepth}{2}
\setcounter{secnumdepth}{3}

\definecolor{DeepBlue}{HTML}{17365D}
\definecolor{SoftBlue}{HTML}{EAF1F8}
\definecolor{DeepGreen}{HTML}{1D5B45}
\definecolor{SoftGreen}{HTML}{EDF7F2}
\definecolor{DeepRed}{HTML}{8A2D2D}
\definecolor{SoftGray}{HTML}{F4F5F7}
\definecolor{LinkBlue}{HTML}{145DA0}

\hypersetup{
  colorlinks=true,
  linkcolor=DeepBlue,
  citecolor=DeepGreen,
  urlcolor=LinkBlue,
  pdftitle={Dyadic Spectral Determinants for the Fabius--Rvachev System},
  pdfauthor={OpenAI GPT-5.6 Pro},
  pdfsubject={Thue--Morse sequence, Fabius function, Rvachev up-function, q-binomial coefficients, and infinite sinc products}
}

\pagestyle{fancy}
\fancyhf{}
\fancyhead[L]{\small\textsc{Dyadic Spectral Determinants}}
\fancyhead[R]{\small\textsc{Fabius--Rvachev Frontier Report}}
\fancyfoot[C]{\thepage}
\renewcommand{\headrulewidth}{0.4pt}

\titleformat{\section}{\Large\bfseries\color{DeepBlue}}{\thesection}{0.8em}{}
\titleformat{\subsection}{\large\bfseries\color{DeepBlue}}{\thesubsection}{0.7em}{}
\titleformat{\subsubsection}{\normalsize\bfseries\color{DeepGreen}}{\thesubsubsection}{0.6em}{}

\theoremstyle{plain}
\newtheorem{theorem}{Theorem}[section]
\newtheorem{proposition}[theorem]{Proposition}
\newtheorem{corollary}[theorem]{Corollary}
\newtheorem{lemma}[theorem]{Lemma}
\newtheorem{conjecture}[theorem]{Conjecture}
\theoremstyle{definition}
\newtheorem{definition}[theorem]{Definition}
\newtheorem{problem}[theorem]{Problem}
\theoremstyle{remark}
\newtheorem{remark}[theorem]{Remark}
\newtheorem{checkitem}[theorem]{Computational check}

\newcommand{\Up}{\operatorname{up}}
\newcommand{\sinc}{\operatorname{sinc}}
\newcommand{\Tr}{\operatorname{Tr}}
\newcommand{\Res}{\operatorname*{Res}}
\newcommand{\FP}{\operatorname*{FP}}
\newcommand{\Li}{\operatorname{Li}}
\newcommand{\arsinh}{\operatorname{arsinh}}
\newcommand{\sgn}{\operatorname{sgn}}
\newcommand{\rank}{\operatorname{rank}}
\newcommand{\dd}{\,\mathrm{d}}
\newcommand{\N}{\mathbb N}
\newcommand{\R}{\mathbb R}
\newcommand{\C}{\mathbb C}
\newcommand{\Z}{\mathbb Z}
\newcommand{\Hh}{\mathcal H}
\newcommand{\Ecal}{\mathscr E}
\newcommand{\Pcal}{\mathcal P}
\newcommand{\Lcal}{\mathcal L}
\newcommand{\Acal}{\mathcal A}
\newcommand{\Dcal}{\mathcal D}
\newcommand{\vtwo}{\nu_2}
\newcommand{\tm}{\tau}
\newcommand{\qpoch}[2]{(#1;#2)_\infty}
\newcommand{\qpochfin}[3]{(#1;#2)_{#3}}
\newcommand{\qbinom}[3]{\genfrac{[}{]}{0pt}{}{#1}{#2}_{#3}}
\newcommand{\e}{\mathrm e}
\newcolumntype{P}[1]{>{\raggedright\arraybackslash}p{#1}}

\newenvironment{statusbox}[1]{%
  \par\medskip\noindent\begin{minipage}{\textwidth}%
  \color{DeepBlue}\hrule\vspace{5pt}%
  \color{black}\textbf{\color{DeepBlue}#1.}\ }{%
  \par\vspace{5pt}\color{DeepBlue}\hrule
  \end{minipage}\par\medskip}

\newenvironment{resultbox}[1]{%
  \par\medskip\noindent\begin{minipage}{\textwidth}%
  \color{DeepGreen}\hrule\vspace{5pt}%
  \color{black}\textbf{\color{DeepGreen}#1.}\ }{%
  \par\vspace{5pt}\color{DeepGreen}\hrule
  \end{minipage}\par\medskip}

\newcommand{\CorpusURL}{\url{https://github.com/VladimirReshetnikov/ProveIt/tree/main/Analysis/FabiusFunction/docs}}

\begin{document}
\pagenumbering{gobble}

\begin{titlepage}
\centering
\vspace*{18mm}
{\Huge\bfseries\color{DeepBlue} Dyadic Spectral Determinants\par}
\vspace{3mm}
{\LARGE\bfseries for the Fabius--Rvachev System\par}
\vspace{7mm}
{\Large New zero-spectrum, local-jet, heat-trace, and $q$-binomial bridges\par}
\vspace{16mm}
\rule{0.82\textwidth}{0.6pt}\par
\vspace{10mm}
{\large A frontier report based on the complete recursive \texttt{*.tex} corpus under\par}
\vspace{3mm}
{\small\ttfamily Analysis/FabiusFunction/docs\par}
\vspace{10mm}
{\large 27 August 2026\par}
\vfill
\begin{minipage}{0.88\textwidth}
\small
\textbf{Central outcome.}
The Fourier transform of Rvachev's up-function is reorganized as the Fredholm determinant of a dyadically self-similar diagonal operator.  This produces, in one framework, an exact spectral zeta function, a complex-dimension heat expansion, exact all-order jets at every integer zero, a multiscale $q$-Euler determinant, finite Gaussian-binomial coefficient formulae, and a symmetric $q$-integer deformation whose $q\uparrow1$ limit is the original infinite sinc product.
\end{minipage}
\vspace{12mm}
\end{titlepage}
\pagenumbering{roman}

\begin{abstract}
Let
\[
 \Phi(z)=\widehat{\Up}(z)
 =\prod_{j=0}^{\infty}\sinc\!\left(\frac{\pi z}{2^j}\right),
 \qquad \sinc w=\frac{\sin w}{w}.
\]
The inspected corpus already contains the canonical integer-zero product, the exact multiplicity $1+\nu_2(m)$ at a nonzero integer $m$, and the Thue--Morse sign law between consecutive zeros.  The present report does not relabel those facts as new.  Instead it introduces the diagonal operator
\[
 D e_{j,n}=2^j n\,e_{j,n}\qquad(j\ge0,\ n\ge1)
\]
and proves
\[
 \Phi(z)=\det(I-z^2D^{-2}),\qquad D\simeq N\oplus2D.
\]
The decomposition gives the exact positive zero count $2M-s_2(M)$, the spectral zeta function
\[
 \zeta_D(s)=\frac{\zeta(s)}{1-2^{-s}},
\]
and a heat trace with an explicit real-analytic log-periodic term.  We derive a normalized, exact all-order Taylor factorization at every integer zero.  We then define a dyadic $q$-Euler determinant
\[
 \Ecal(u;Q)=\det(I-uQ^D)
 =\prod_{j\ge0}(uQ^{2^j};Q^{2^j})_\infty,
\]
obtain finite multiscale Gaussian-binomial expansions, and calculate the complete small-$t$ Mellin asymptotics of $\Ecal(u;\e^{-t})$, including the critical $u=1$ quadratic logarithm.  Finally, with the symmetric $q$-integer $[m]_Q$, we construct
\[
 \Phi_Q(z)=\prod_{m\ge1}\left(1-\frac{z^2}{[m]_Q^2}\right)^{1+\nu_2(m)}
\]
and prove an exact $q$-Pochhammer factorization, the scaling law
\[
 \Phi_Q(z)=S_Q(z)\Phi_{Q^2}\!\left(z/[2]_Q\right),
\]
and locally uniform convergence $\Phi_Q\to\Phi$ as $Q\uparrow1$.  This supplies a direct structural bridge from the Fourier zero divisor to $q$-Pochhammer symbols and $q$-binomial coefficients.
\end{abstract}

\tableofcontents
\clearpage
\pagenumbering{arabic}

\section{Scope, corpus audit, and novelty protocol}

\subsection{What was inspected}

The recursive documentation tree at \CorpusURL{} is not a single article.  It is a stratified mathematical record containing living expositions, an explicitly frozen historical archive, a consolidated non-formalized frontier volume, a Thue--Morse formula atlas, twelve successive Fourier-decay analyses, and imported source transcriptions of external papers.  The corpus ledger in \cref{app:ledger} records $57$ TeX artifacts:
\begin{center}
\begin{tabular}{@{}lrr@{}}
\toprule
Stratum & TeX artifacts & Role\\
\midrule
Living and canonical documents & 18 & Current expositions and active frontier drafts\\
Frozen historical archive & 32 & Provenance, superseded drafts, and deltas\\
Imported paper sources & 7 & External mathematical context\\
\midrule
Total & 57 & Recursive corpus audited\\
\bottomrule
\end{tabular}
\end{center}

Every TeX artifact was included in the audit.  Mathematical extraction was concentrated on the living/canonical descendants and the active Fourier series; frozen ancestors were compared through their archived first-version source, provenance descriptions, and surviving distinctive sections.  This distinction matters: the archive intentionally preserves superseded states and explicitly prunes byte-identical copies.  Treating all $57$ files as independent contemporary claims would create false novelty and heavy duplication.

\subsection{Three status classes}

\begin{statusbox}{Status convention}
Throughout this report, \emph{established baseline} means a statement already present in the inspected repository or in a cited external source.  \emph{New deduction} means that the theorem and proof were not found in the inspected corpus or in targeted searches made for this report.  This is a relative corpus claim, not a universal priority claim.  The proofs below are mathematical proofs written for human verification.  Except where an explicit Lean-status note cites a declaration, they have not yet been checked by Lean.
\end{statusbox}

The main examples of statements that are \emph{not} claimed here as new are:
\begin{enumerate}[label=(\roman*),leftmargin=2.1em]
\item the infinite sinc product for $\widehat{\Up}$ and its dyadic scaling equation;
\item the canonical product
\[
 \Phi(z)=\prod_{m=1}^{\infty}\left(1-\frac{z^2}{m^2}\right)^{1+\nu_2(m)};
\]
\item the assertion that the nonzero zeros are precisely the integers, with multiplicity $1+\nu_2(|m|)$;
\item strict logarithmic concavity of $|\Phi|$ between consecutive zeros and the sign formula
\[
 \sgn\Phi(x)=(-1)^{s_2(M)}\qquad(M<x<M+1);
\]
\item the existing moment/cumulant formulae for the Rvachev probability law;
\item the base-$1/2$ $q$-Pochhammer and Gaussian-binomial formulae for exact dyadic values of the Fabius function;
\item the pointwise, metric, mean, extremal, and annular Fourier-decay results assembled in the final Fourier-decay synthesis.
\end{enumerate}

The purpose of the report is instead to place those facts inside a new spectral determinant and $q$-deformation architecture and to extract consequences not present in the source corpus.

\subsection{Result map}

\begin{resultbox}{New deductions developed below}
The report proves: (a) a Fredholm determinant model with $D\simeq N\oplus2D$; (b) exact zero-counting and spectral-parity formulae; (c) the meromorphic spectral zeta $\zeta(s)/(1-2^{-s})$ and its heat trace with explicit log-periodic oscillation; (d) exact normalized all-order jets at every integer zero; (e) an operator-valued dyadic $q$-Euler product and finite multiscale Gaussian-binomial expansion; (f) complete small-$t$ asymptotics of the resulting tower of $q$-Pochhammer symbols; (g) a symmetric $q$-integer deformation of $\Phi$ with exact $q$-scaling, zero arithmetic, and $q\uparrow1$ convergence; and (h) a base-$b$ extension explaining which features are genuinely dyadic.
\end{resultbox}

\section{Established baseline and notation}

\subsection{The twin functions}

Rvachev's up-function $\Up$ is the even $C^\infty$ density supported on $[-1,1]$ whose Fourier transform, under the convention
\[
 \widehat f(z)=\int_{\R}f(x)\e^{-2\pi i z x}\dd x,
\]
is
\begin{equation}\label{eq:Phi-sinc}
 \Phi(z):=\widehat{\Up}(z)
 =\prod_{j=0}^{\infty}\sinc\!\left(\frac{\pi z}{2^j}\right).
\end{equation}
The Fabius function $F$ is its affine cumulative twin; on the support one may write
\[
 \Up(x)=
 \begin{cases}
 F(x+1),&-1\le x\le0,\\
 F(1-x),&0\le x\le1.
 \end{cases}
\]
For the present report the Fourier entire function $\Phi$ is the central object, while the relation to $F$ explains why the resulting identities belong to the same theory.

A probability realization is useful.  If $U_0,U_1,\ldots$ are independent and uniform on $[-1/2,1/2]$, then
\begin{equation}\label{eq:prob-sum}
 X=\sum_{j=0}^{\infty}2^{-j}U_j
\end{equation}
has density $\Up$ and characteristic function $\mathbb E\e^{itX}=\Phi(t/(2\pi))$.

\subsection{Thue--Morse and binary notation}

For $n\ge0$, let $s_2(n)$ denote the sum of the binary digits of $n$, and put
\begin{equation}\label{eq:tm}
 \tm_n=(-1)^{s_2(n)}.
\end{equation}
The sequence satisfies $\tm_{2n}=\tm_n$ and $\tm_{2n+1}=-\tm_n$.  We write $\nu_2(m)$ for the exponent of $2$ in a positive integer $m$.

\subsection{$q$-notation}

For $|Q|<1$,
\[
 (a;Q)_L=\prod_{r=0}^{L-1}(1-aQ^r),
 \qquad
 (a;Q)_\infty=\prod_{r=0}^{\infty}(1-aQ^r),
\]
and the Gaussian binomial coefficient is
\[
 \qbinom{L}{k}{Q}=\frac{(Q;Q)_L}{(Q;Q)_k(Q;Q)_{L-k}}.
\]
The finite $q$-binomial theorem is
\begin{equation}\label{eq:q-binomial}
 (a;Q)_L=\sum_{k=0}^{L}(-1)^kQ^{\binom{k}{2}}
 \qbinom{L}{k}{Q}a^k.
\end{equation}
The corpus uses these functions chiefly at base $Q=1/2$ to encode exact dyadic values of $F$.  The constructions below use a different $Q$: it deforms the \emph{Fourier zero spectrum}.  The two roles should not be conflated.

\subsection{The established zero divisor}

Euler's sine product applied scale by scale to \cref{eq:Phi-sinc} gives
\begin{equation}\label{eq:canonical-established}
 \Phi(z)=\prod_{j=0}^{\infty}\prod_{n=1}^{\infty}
 \left(1-\frac{z^2}{(2^jn)^2}\right)
 =\prod_{m=1}^{\infty}
 \left(1-\frac{z^2}{m^2}\right)^{a_m},
 \qquad a_m:=1+\nu_2(m).
\end{equation}
The second equality counts the representations $m=2^jn$.  Consequently the nonzero zeros are exactly $\pm m$, each with multiplicity $a_m$.  The active Fourier-decay corpus already contains \cref{eq:canonical-established}; it is the starting point, not a new theorem of this report.

\section{The dyadic spectral operator}

\subsection{Definition and self-similarity}

Let
\[
 \Hh=\bigoplus_{j=0}^{\infty}\ell^2(\N),
\]
with orthonormal basis $e_{j,n}$, $j\ge0$, $n\ge1$.  Define the positive diagonal operator
\begin{equation}\label{eq:D-def}
 De_{j,n}=2^jn\,e_{j,n}.
\end{equation}
Let $N$ be the number operator on $\ell^2(\N)$, $Ne_n=ne_n$.

\begin{theorem}[Dyadic spectral decomposition]\label{thm:D-decomp}
There is a unitary equivalence
\begin{equation}\label{eq:D-decomp}
 D\simeq N\oplus2D.
\end{equation}
Moreover $D^{-p}$ is trace class precisely for $p>1$.
\end{theorem}

\begin{proof}
The subspace with $j=0$ carries the eigenvalues $n$, hence is $N$.  Relabeling $e_{j,n}$ as $e_{j-1,n}$ on the subspace $j\ge1$ identifies the restriction with $2D$.  Finally,
\[
 \Tr D^{-p}
 =\sum_{j=0}^{\infty}\sum_{n=1}^{\infty}(2^jn)^{-p}
 =\frac{\zeta(p)}{1-2^{-p}},
\]
which is finite exactly when $p>1$.
\end{proof}

The decomposition is the operator-theoretic form of dyadic refinement.  It is also the source of every denominator $1-2^{-s}$ below.

\subsection{The sinc product as a Fredholm determinant}

Since $D^{-2}$ is trace class, the ordinary Fredholm determinant $\det(I-z^2D^{-2})$ is defined for every $z\in\C$.

\begin{theorem}[Fredholm determinant representation]\label{thm:fredholm}
For every $z\in\C$,
\begin{equation}\label{eq:fredholm}
 \boxed{\ \Phi(z)=\det(I-z^2D^{-2}).\ }
\end{equation}
The operator decomposition \cref{eq:D-decomp} implies the exact scaling equation
\begin{equation}\label{eq:scaling-spectral}
 \Phi(z)=\sinc(\pi z)\Phi(z/2).
\end{equation}
\end{theorem}

\begin{proof}
The eigenvalues of $D^{-2}$ are $(2^jn)^{-2}$, and their sum is finite.  Thus
\[
 \det(I-z^2D^{-2})
 =\prod_{j=0}^{\infty}\prod_{n=1}^{\infty}
 \left(1-\frac{z^2}{(2^jn)^2}\right).
\]
Euler's product $\sinc(\pi w)=\prod_{n\ge1}(1-w^2/n^2)$ identifies this with \cref{eq:Phi-sinc}.  Applying the multiplicativity of the determinant to $D\simeq N\oplus2D$ gives
\[
 \det(I-z^2D^{-2})
 =\det(I-z^2N^{-2})
  \det\!\left(I-z^2(2D)^{-2}\right)
 =\sinc(\pi z)\Phi(z/2).
\]
\end{proof}

\begin{remark}
This is an \emph{ordinary} Fredholm determinant, not a zeta-regularized determinant.  The trace-class input is $D^{-2}$, and no renormalization is required.
\end{remark}

\subsection{Multiplicity as spectral degeneracy}

The eigenvalue $m$ of $D$ occurs once for every $j$ such that $2^j\mid m$.  Therefore
\begin{equation}\label{eq:am}
 \operatorname{mult}_D(m)=a_m=1+\nu_2(m),
 \qquad
 a_{2m}=a_m+1,
 \qquad
 a_{2m+1}=1.
\end{equation}
Thus the canonical zero multiplicity is exactly the spectral degeneracy of $D$.  The elementary $2$-regular recursion in \cref{eq:am} is the discrete shadow of $D\simeq N\oplus2D$.

\section{Exact zero arithmetic and spectral parity}

\subsection{The positive counting function}

Let
\[
 A(M)=\rank\mathbf 1_{[1,M]}(D)
 =\sum_{m=1}^{M}a_m
\]
count positive eigenvalues, including multiplicity.

\begin{theorem}[Exact Weyl law with binary remainder]\label{thm:counting}
For every integer $M\ge0$,
\begin{equation}\label{eq:A-exact}
 \boxed{\ A(M)=\sum_{j=0}^{\infty}\left\lfloor\frac{M}{2^j}\right\rfloor
 =2M-s_2(M).\ }
\end{equation}
Consequently, with $M=\lfloor R\rfloor$, the number of nonzero zeros of $\Phi$ in $|z|\le R$, counted with multiplicity, is
\begin{equation}\label{eq:Nphi}
 \boxed{\ \mathcal N_\Phi(R)=4M-2s_2(M).\ }
\end{equation}
\end{theorem}

\begin{proof}
Counting the pairs $(j,n)$ with $2^jn\le M$ gives the floor sum.  On the other hand,
\[
 A(M)=M+\sum_{m=1}^{M}\nu_2(m)
 =M+\nu_2(M!).
\]
Legendre's binary formula $\nu_2(M!)=M-s_2(M)$ proves \cref{eq:A-exact}.  The zeros are symmetric, so the disk count is twice $A(M)$.
\end{proof}

The exact formula refines the linear Weyl term $A(M)\sim2M$ by a digit-sum fluctuation.  In particular,
\begin{equation}\label{eq:A-bounds}
 2M-\lfloor\log_2M\rfloor-1\le A(M)\le2M-1
 \qquad(M\ge1).
\end{equation}
The recursions
\begin{equation}\label{eq:A-rec}
 A(2M)=A(M)+2M,
 \qquad
 A(2M+1)=A(M)+2M+1
\end{equation}
follow immediately.

\subsection{Thue--Morse as spectral-flow parity}

\begin{corollary}[Spectral parity formula]\label{cor:spectral-parity}
For every $M\ge0$,
\begin{equation}\label{eq:spectral-parity}
 \boxed{\ (-1)^{\rank\mathbf 1_{[1,M]}(D)}=(-1)^{A(M)}
 =(-1)^{s_2(M)}=\tm_M.\ }
\end{equation}
Hence, for $M<x<M+1$,
\[
 \sgn\Phi(x)=(-1)^{\rank\mathbf 1_{[1,M]}(D)}=\tm_M.
\]
\end{corollary}

\begin{proof}
The even term $2M$ disappears modulo $2$ in \cref{eq:A-exact}.  Starting with $\Phi(0)=1$, crossing the zeros $1,\ldots,M$ multiplies the sign by $(-1)^{A(M)}$.
\end{proof}

The interval-sign theorem was already present in the Fourier corpus.  What is new here is its interpretation: the Thue--Morse bit is the parity of a finite-rank spectral projection of a self-similar operator.

\subsection{Dyadic-block statistics}

For uniform $M$ in $\{0,\ldots,2^n-1\}$, the random variable $s_2(M)$ is binomial $\operatorname{Bin}(n,1/2)$.  Thus
\begin{equation}\label{eq:block-mean}
 \frac1{2^n}\sum_{M=0}^{2^n-1}\bigl(2M-A(M)\bigr)=\frac n2,
 \qquad
 \operatorname{Var}(2M-A(M))=\frac n4.
\end{equation}
The de Moivre--Laplace theorem gives
\begin{equation}\label{eq:count-clt}
 \frac{A(M)-(2M-n/2)}{\sqrt n/2}
 \Longrightarrow \mathcal N(0,1)
\end{equation}
under dyadic-block sampling.  Thus the deterministic zero-count remainder has an exact probabilistic law when sampled on its natural self-similar blocks.

\subsection{A certified canonical-product truncation}

Let
\[
 \Phi_M(z)=\prod_{m=1}^{M}\left(1-\frac{z^2}{m^2}\right)^{a_m}.
\]

\begin{proposition}[Elementary tail bound]\label{prop:tail}
For $M\ge1$,
\begin{equation}\label{eq:trace-tail}
 T_M:=\sum_{m>M}\frac{a_m}{m^2}
 <\frac2M+\frac{2\zeta(2)}{M^2}.
\end{equation}
If $|z|\le R$ and $M\ge\sqrt2R$, then, on any zero-free simply connected region where the logarithm is chosen continuously,
\begin{equation}\label{eq:det-tail}
 \left|\log\frac{\Phi(z)}{\Phi_M(z)}\right|
 \le2R^2T_M.
\end{equation}
\end{proposition}

\begin{proof}
Write the tail spectrally:
\[
 T_M=\sum_{j=0}^{\infty}4^{-j}
 \sum_{n>M/2^j}\frac1{n^2}.
\]
For $j\le\lfloor\log_2M\rfloor$, the inner tail is at most $2^j/M$; summing gives less than $2/M$.  For the remaining $j$, bound the inner sum by $\zeta(2)$ and the geometric tail by $2/M^2$.  If $|z|^2/m^2\le1/2$, then $|\log(1-w)|\le2|w|$, and summation proves \cref{eq:det-tail}.
\end{proof}

\section{Spectral zeta, determinant coefficients, and cumulants}

\subsection{The exact spectral zeta function}

\begin{theorem}[Dyadic spectral zeta]\label{thm:zetaD}
For $\Re s>1$,
\begin{equation}\label{eq:zetaD}
 \boxed{\ \zeta_D(s):=\Tr D^{-s}
 =\sum_{m=1}^{\infty}\frac{1+\nu_2(m)}{m^s}
 =\frac{\zeta(s)}{1-2^{-s}}.\ }
\end{equation}
This identity gives a meromorphic continuation to $\C$.  Besides the pole at $s=1$ with residue $2$, it has the dyadic pole lattice
\begin{equation}\label{eq:omega}
 \omega_k=\frac{2\pi i k}{\log2},\qquad k\in\Z,
\end{equation}
with
\begin{equation}\label{eq:zetaD-res}
 \Res_{s=\omega_k}\zeta_D(s)=\frac{\zeta(\omega_k)}{\log2}.
\end{equation}
In particular,
\begin{equation}\label{eq:zetaD-zeroexp}
 \zeta_D(s)=-\frac{1}{2\log2}\frac1s
 -\frac14-\frac{\log(2\pi)}{2\log2}+O(s)
 \qquad(s\to0).
\end{equation}
\end{theorem}

\begin{proof}
The double sum in the proof of \cref{thm:D-decomp} gives \cref{eq:zetaD}.  The pole at $1$ is multiplied by $(1-2^{-1})^{-1}=2$.  At a zero $\omega_k$ of $1-2^{-s}$, its derivative is $(\log2)2^{-s}=\log2$, which gives \cref{eq:zetaD-res}.  Expansion of $\zeta(s)$ and $(1-\e^{-s\log2})^{-1}$ at $0$ gives \cref{eq:zetaD-zeroexp}.
\end{proof}

For $k\ne0$, the numerator does not vanish on the imaginary axis; this follows, for example, from the functional equation together with the classical zero-free line $\Re s=1$.  The points $\omega_k$ therefore form a genuine vertical lattice of complex dimensions.

\begin{corollary}[Reciprocal zero power sums]\label{cor:zero-powers}
For every integer $r\ge1$,
\begin{equation}\label{eq:zero-powers}
 \sum_{m=1}^{\infty}\frac{1+\nu_2(m)}{m^{2r}}
 =\frac{\zeta(2r)}{1-4^{-r}}.
\end{equation}
\end{corollary}

\subsection{Trace logarithm and a coefficient recurrence}

For $|z|<1$, the trace logarithm of the Fredholm determinant is absolutely convergent.

\begin{theorem}[Spectral Taylor expansion]\label{thm:taylor}
One has
\begin{equation}\label{eq:logPhi}
 \boxed{\ \log\Phi(z)
 =-\sum_{r=1}^{\infty}
 \frac{\zeta(2r)}{r(1-4^{-r})}z^{2r},
 \qquad |z|<1.\ }
\end{equation}
If
\[
 \Phi(z)=\sum_{n=0}^{\infty}c_nz^{2n},\qquad c_0=1,
\]
then
\begin{equation}\label{eq:c-rec}
 \boxed{\ c_n=-\frac1n\sum_{r=1}^{n}
 \frac{\zeta(2r)}{1-4^{-r}}c_{n-r}.\ }
\end{equation}
In particular,
\begin{equation}\label{eq:first-c}
 c_1=-\frac{2\pi^2}{9},
 \qquad
 c_2=\frac{38\pi^4}{2025}.
\end{equation}
\end{theorem}

\begin{proof}
For $|z|<1$,
\[
 \log\det(I-z^2D^{-2})
 =-\sum_{r\ge1}\frac{z^{2r}}r\Tr D^{-2r},
\]
and \cref{cor:zero-powers} gives \cref{eq:logPhi}.  Put $x=z^2$ and differentiate logarithmically:
\[
 x\frac{\sum_{n\ge1}nc_nx^{n-1}}{\sum_{n\ge0}c_nx^n}
 =-\sum_{r\ge1}\frac{\zeta(2r)}{1-4^{-r}}x^r.
\]
Coefficient comparison gives \cref{eq:c-rec}.
\end{proof}

\subsection{A spectral interpretation of the known cumulants}

Let $X$ be the centered random variable in \cref{eq:prob-sum}.  Since $\mathbb E\e^{itX}=\Phi(t/(2\pi))$, comparison of \cref{eq:logPhi} with the cumulant expansion gives
\begin{equation}\label{eq:cumulants}
 \kappa_{2r}(X)=\frac{B_{2r}}{2r(1-4^{-r})},
 \qquad
 \kappa_{2r+1}(X)=0.
\end{equation}
These cumulants already occur in the frontier corpus.  The new point is the identity
\begin{equation}\label{eq:cumulant-trace}
 \kappa_{2r}(X)
 =(-1)^{r+1}\frac{(2r)!}{r(2\pi)^{2r}}\Tr D^{-2r},
\end{equation}
which interprets them as spectral power traces.  For example $\kappa_2=1/9$ follows from $\Tr D^{-2}=2\pi^2/9$.

\section{Heat trace and dyadic complex dimensions}

\subsection{Mellin representation}

Define
\begin{equation}\label{eq:Theta}
 \Theta(t)=\Tr\e^{-tD}
 =\sum_{m=1}^{\infty}a_m\e^{-tm}
 =\sum_{j=0}^{\infty}\frac1{\e^{2^jt}-1},
 \qquad t>0.
\end{equation}
For $\Re s>1$,
\begin{equation}\label{eq:Theta-Mellin}
 \int_0^{\infty}\Theta(t)t^{s-1}\dd t
 =\Gamma(s)\zeta_D(s)
 =\Gamma(s)\frac{\zeta(s)}{1-2^{-s}}.
\end{equation}
The pole lattice \cref{eq:omega} therefore produces logarithmic periodicity in physical scale.

Put
\[
 L=\log2,
 \qquad
 \omega_k=\frac{2\pi i k}{L},
\]
and define the real-analytic, one-periodic function
\begin{equation}\label{eq:Ptheta}
 \Pcal_\Theta(u)=\frac1L\sum_{k\in\Z\setminus\{0\}}
 \Gamma(\omega_k)\zeta(\omega_k)\e^{2\pi iku}.
\end{equation}
The exponential decay of $\Gamma(i y)$ makes the Fourier series absolutely and locally uniformly convergent with all derivatives.

\subsection{Complete heat expansion}

\begin{theorem}[Heat trace with log-periodic oscillation]\label{thm:heat}
Let
\begin{equation}\label{eq:Ctheta}
 C_\Theta=\frac{\gamma-\log(2\pi)}{2\log2}-\frac14.
\end{equation}
For every integer $J\ge0$ and every $0<\varepsilon<1$,
\begin{align}\label{eq:heat-asymp}
 \Theta(t)
 ={}&\frac2t-\frac{\log(1/t)}{2\log2}+C_\Theta
 +\Pcal_\Theta\!\left(\log_2\frac1t\right)\\
 &-\sum_{r=1}^{J}
 \frac{B_{2r}}{(2r)!(2^{2r-1}-1)}t^{2r-1}
 +O_{J,\varepsilon}\!\left(t^{2J+1-\varepsilon}\right)
 \qquad(t\downarrow0).\notag
\end{align}
The first algebraic correction is $-t/12$.
\end{theorem}

\begin{proof}
Mellin inversion gives
\[
 \Theta(t)=\frac1{2\pi i}\int_{c-i\infty}^{c+i\infty}
 \Gamma(s)\zeta_D(s)t^{-s}\dd s,
 \qquad c>1.
\]
Move the contour left.  The pole at $s=1$ contributes $2/t$.  At $s=0$, $\Gamma(s)\zeta_D(s)$ has a double pole.  From \cref{eq:zetaD-zeroexp},
\[
 \Gamma(s)\zeta_D(s)
 =-\frac1{2L}\frac1{s^2}
 +\left(\frac{\gamma-\log(2\pi)}{2L}-\frac14\right)\frac1s+O(1).
\]
Multiplication by $t^{-s}$ gives the logarithm and $C_\Theta$.  The residues at $s=\omega_k$, $k\ne0$, sum to \cref{eq:Ptheta}.  At the negative odd integer $s=-(2r-1)$, the residue of $\Gamma$ is $-1/(2r-1)!$ and $\zeta(1-2r)=-B_{2r}/(2r)$; this gives the displayed coefficient.  The trivial zeros of $\zeta$ cancel the poles of $\Gamma$ at negative even integers.  Standard vertical-line bounds and the exponential decay of $\Gamma$ justify the contour shift and remainder.
\end{proof}

\subsection{Size of the oscillation}

The first Fourier coefficient in \cref{eq:Ptheta} is
\begin{equation}\label{eq:Pfirst}
 \frac{\Gamma(2\pi i/L)\zeta(2\pi i/L)}L
 =-8.226408112794756\times10^{-7}
  -9.760012263161017\times10^{-7}i,
\end{equation}
with modulus $1.276446747088608\times10^{-6}$.  Thus the oscillation is small but exact.  It is not an error term: it is a Fourier series forced by the pole lattice of $\zeta_D$.

\begin{remark}[Two kinds of self-similar fluctuation]
The exact counting remainder $-s_2(M)$ is discontinuous and arithmetic.  The heat remainder $\Pcal_\Theta(\log_2(1/t))$ is smooth and analytic.  They are two transforms of the same dyadic spectral multiplicity: a sharp cutoff retains binary carries, whereas exponential smoothing turns the scale lattice into an analytic periodic wave.
\end{remark}

\section{Exact local jets at all integer zeros}

\subsection{A normalized factorization}

Fix a nonzero integer $m$.  Write
\begin{equation}\label{eq:m-rq}
 r=\nu_2(|m|),\qquad a=r+1,\qquad q=\frac{m}{2^r},
\end{equation}
so that $q$ is odd.

\begin{theorem}[Exact all-order zero jet]\label{thm:local-jet}
For $w$ in a neighborhood of $0$,
\begin{equation}\label{eq:local-factor}
 \Phi(m+w)=C_mw^aR_m(w),
\end{equation}
where
\begin{equation}\label{eq:Cm}
 \boxed{\ C_m=-\frac{\Phi(q/2)}{m^a}\ }
\end{equation}
and
\begin{equation}\label{eq:Rm}
 R_m(w)=
 \left(\frac{m}{m+w}\right)^a
 \prod_{h=0}^{r}\sinc\!\left(\frac{\pi w}{2^h}\right)
 \frac{\Phi\!\left(q/2+w/2^{r+1}\right)}{\Phi(q/2)}.
\end{equation}
In particular $R_m$ is analytic and $R_m(0)=1$, and every derivative is given exactly by
\begin{equation}\label{eq:all-jets}
 \boxed{\ \Phi^{(a+\ell)}(m)
 =C_m(a+\ell)![w^\ell]R_m(w),\qquad \ell\ge0.\ }
\end{equation}
The first nonzero derivative is
\begin{equation}\label{eq:first-derivative-zero}
 \boxed{\ \Phi^{(a)}(m)=-\frac{a!}{m^a}\Phi(q/2).\ }
\end{equation}
\end{theorem}

\begin{proof}
For $0\le h\le r$, $m/2^h$ is an integer and
\begin{align*}
 \sinc\!\left(\frac{\pi(m+w)}{2^h}\right)
 &=(-1)^{m/2^h}\frac{w}{m+w}
   \sinc\!\left(\frac{\pi w}{2^h}\right).
\end{align*}
The product of the signs is
\[
 (-1)^{\sum_{h=0}^{r}m/2^h}
 =(-1)^{q(2^{r+1}-1)}=-1.
\]
The tail of the original sinc product is
\[
 \prod_{h=r+1}^{\infty}
 \sinc\!\left(\frac{\pi(m+w)}{2^h}\right)
 =\Phi\!\left(\frac q2+\frac{w}{2^{r+1}}\right).
\]
Multiplying these identities gives
\[
 \Phi(m+w)=-\left(\frac{w}{m+w}\right)^a
 \prod_{h=0}^{r}\sinc\!\left(\frac{\pi w}{2^h}\right)
 \Phi\!\left(\frac q2+\frac{w}{2^{r+1}}\right),
\]
which is equivalent to \cref{eq:local-factor}--\cref{eq:Rm}.  Since $q/2$ is not an integer, $\Phi(q/2)\ne0$.  Coefficient extraction proves \cref{eq:all-jets}.
\end{proof}

\begin{statusbox}{Lean status of the local-jet theorem}
The denominator-cleared algebraic factorization is now machine checked,
globally in the complex local coordinate and for every nonzero integer
center, as
\path{Fabius.rvachevFourierProduct_int_add_factorization} in
\path{IntegerZeroLocalFactorization.lean}.  The divided equality holds as a
germ at zero by
\path{Fabius.rvachevFourierProduct_int_add_eventuallyEq_pow_mul_cofactor}.
The resulting cofactor is analytic and nonzero at the origin; Lean proves
its exact value, the exact analytic order \(1+\nu_2(|m|)\), every lower
derivative vanishing, and
\[
 \Phi^{(a)}(m)=-\frac{a!}{m^a}\Phi(q/2)
\]
through the declarations in \path{IntegerZeroAnalyticOrder.lean}.  The
field-generic extraction of this leading derivative is in
\path{LeadingJet.lean}.  Thus
\cref{eq:local-factor,eq:Cm,eq:first-derivative-zero} and the
\(\ell=0\) case of \cref{eq:all-jets} are formalized.  The explicitly
normalized function \(R_m\), the coefficient formula
\cref{eq:all-jets} for \(\ell>0\), the next-jet formula, dyadic-ray
recursion, and the packaged leading-sign corollary do not yet have matching
named Lean declarations.
\end{statusbox}

\subsection{The next coefficient and dyadic-ray recursion}

Differentiating \cref{eq:Rm} at $0$ gives
\begin{equation}\label{eq:Rprime}
 R_m'(0)=-\frac{a}{m}
 +\frac1{2^{r+1}}\frac{\Phi'(q/2)}{\Phi(q/2)}.
\end{equation}
Therefore
\begin{equation}\label{eq:next-jet}
 \Phi^{(a+1)}(m)
 =-\frac{(a+1)!\Phi(q/2)}{m^a}
 \left(-\frac{a}{m}
 +\frac1{2^{r+1}}\frac{\Phi'(q/2)}{\Phi(q/2)}\right).
\end{equation}
All higher jets follow by formal multiplication of the three analytic factors in \cref{eq:Rm}.

The leading coefficients propagate exactly along dyadic rays.

\begin{corollary}[Dyadic-ray recursion]\label{cor:dyadic-ray}
If $a=1+\nu_2(m)$, then
\begin{equation}\label{eq:C-ray}
 C_{2m}=\frac{C_m}{2^{a+1}m}.
\end{equation}
Equivalently,
\begin{equation}\label{eq:der-ray}
 \frac{\Phi^{(a+1)}(2m)}{(a+1)!}
 =\frac1{2^{a+1}m}\frac{\Phi^{(a)}(m)}{a!}.
\end{equation}
\end{corollary}

\begin{proof}
The integers $m$ and $2m$ have the same odd part, while their multiplicities differ by one.  Substitute into \cref{eq:Cm}.
\end{proof}

\subsection{Thue--Morse signs of the leading jets}

For $m>0$, write its odd part as $q=2k+1$.  The established interval-sign theorem gives $\sgn\Phi(q/2)=(-1)^{s_2(k)}$.  Hence
\begin{equation}\label{eq:jet-sign}
 \boxed{\ \sgn\Phi^{(1+\nu_2(m))}(m)
 =-(-1)^{s_2((q-1)/2)}.\ }
\end{equation}
Thus the leading Taylor coefficient at every positive zero is controlled by the Thue--Morse bit of half the predecessor of its odd part.

\subsection{Numerical checks}

The following values were computed independently from a high-precision truncated sinc product and from \cref{eq:first-derivative-zero}; the two evaluations agreed to at least $70$ decimal digits before rounding.
\begin{center}
\begin{tabular}{@{}rrrr@{}}
\toprule
$m$ & $a=1+\nu_2(m)$ & $\Phi^{(a)}(m)$ & sign\\
\midrule
1  & 1 & $-0.5537712758888100953$ & $-$\\
2  & 2 & $-0.2768856379444050477$ & $-$\\
3  & 1 & $\phantom{-}0.01540076637810223766$ & $+$\\
4  & 3 & $-0.05191605711457594644$ & $-$\\
6  & 2 & $\phantom{-}0.002566794396350372944$ & $+$\\
8  & 4 & $-0.003244753569660996652$ & $-$\\
12 & 3 & $\phantom{-}0.0001604246497718983090$ & $+$\\
\bottomrule
\end{tabular}
\end{center}
These computations are checks, not substitutes for the proof.

\section{The dyadic $q$-Euler determinant}

\subsection{Definition and $q$-Pochhammer tower}

Fix $0<Q<1$.  Since $Q^D$ is trace class, define
\begin{equation}\label{eq:E-def}
 \Ecal(u;Q)=\det(I-uQ^D).
\end{equation}

\begin{theorem}[Dyadic $q$-Euler product]\label{thm:E-product}
For every $u\in\C$,
\begin{align}
 \Ecal(u;Q)
 &=\prod_{m=1}^{\infty}(1-uQ^m)^{a_m}\label{eq:E-mult}\\
 &=\prod_{j=0}^{\infty}(uQ^{2^j};Q^{2^j})_\infty.\label{eq:E-tower}
\end{align}
It satisfies the exact renormalization equation
\begin{equation}\label{eq:E-scale}
 \boxed{\ \Ecal(u;Q)=(uQ;Q)_\infty\Ecal(u;Q^2).\ }
\end{equation}
\end{theorem}

\begin{proof}
The first product is the eigenvalue formula for a Fredholm determinant.  The decomposition $D\simeq N\oplus2D$ gives
\[
 \det(I-uQ^D)=\det(I-uQ^N)\det(I-uQ^{2D})
 =(uQ;Q)_\infty\Ecal(u;Q^2).
\]
Iteration gives \cref{eq:E-tower}; convergence follows from $\sum_{j,n}|u|Q^{2^jn}<\infty$.
\end{proof}

\subsection{Lambert series and plethystic logarithm}

For $|u|Q<1$,
\begin{equation}\label{eq:E-log}
 -\log\Ecal(u;Q)
 =\sum_{\ell=1}^{\infty}\frac{u^\ell}{\ell}
 \sum_{j=0}^{\infty}\frac{Q^{2^j\ell}}{1-Q^{2^j\ell}}.
\end{equation}
Equivalently, the ordinary generating function of the zero multiplicities is the dyadic Lambert series
\begin{equation}\label{eq:a-ogf}
 \boxed{\ \sum_{m=1}^{\infty}a_mx^m
 =\sum_{j=0}^{\infty}\frac{x^{2^j}}{1-x^{2^j}},\qquad |x|<1.\ }
\end{equation}
The same self-similarity now appears in three transforms:
\[
 a_m\quad\longleftrightarrow\quad
 \frac{\zeta(s)}{1-2^{-s}}
 \quad\longleftrightarrow\quad
 \sum_{j\ge0}\frac{x^{2^j}}{1-x^{2^j}}.
\]

\section{Finite multiscale Gaussian-binomial formulae}

\subsection{A finite determinant}

For $M\ge1$, define
\begin{equation}\label{eq:EM}
 \Ecal_M(u;Q)=\prod_{m=1}^{M}(1-uQ^m)^{a_m}.
\end{equation}
Put
\[
 J=\lfloor\log_2M\rfloor,
 \qquad
 L_j=\left\lfloor\frac{M}{2^j}\right\rfloor.
\]

\begin{theorem}[Finite multiscale factorization]\label{thm:finite-E}
One has
\begin{equation}\label{eq:finite-pch}
 \boxed{\ \Ecal_M(u;Q)
 =\prod_{j=0}^{J}(uQ^{2^j};Q^{2^j})_{L_j}.\ }
\end{equation}
Its degree in $u$ is
\begin{equation}\label{eq:EM-degree}
 \deg_u\Ecal_M=A(M)=2M-s_2(M).
\end{equation}
Writing
\begin{equation}\label{eq:EM-coeff-def}
 \Ecal_M(u;Q)=\sum_{k=0}^{A(M)}(-1)^kC_{M,k}(Q)u^k,
\end{equation}
one has the explicit Gaussian-binomial convolution
\begin{equation}\label{eq:CMk}
 \boxed{
 C_{M,k}(Q)=
 \sum_{\substack{k_0+\cdots+k_J=k\\0\le k_j\le L_j}}
 Q^{\sum_{j=0}^{J}2^j\binom{k_j+1}{2}}
 \prod_{j=0}^{J}\qbinom{L_j}{k_j}{Q^{2^j}}.\ }
\end{equation}
\end{theorem}

\begin{proof}
Expand $a_m$ as the number of representations $m=2^jn$ and regroup:
\[
 \Ecal_M(u;Q)
 =\prod_{j=0}^{J}\prod_{n=1}^{L_j}(1-uQ^{2^jn}),
\]
which is \cref{eq:finite-pch}.  Its degree is $\sum_jL_j=A(M)$.  Applying \cref{eq:q-binomial} to each factor gives
\[
 (uQ^{2^j};Q^{2^j})_{L_j}
 =\sum_{k_j=0}^{L_j}(-1)^{k_j}
 Q^{2^j\binom{k_j+1}{2}}
 \qbinom{L_j}{k_j}{Q^{2^j}}u^{k_j}.
\]
Multiplication and collection of the total degree proves \cref{eq:CMk}.
\end{proof}

\subsection{Dyadic endpoints and a multiscale factorial}

For $M=2^n-1$,
\begin{equation}\label{eq:dyadic-L}
 L_j=2^{n-j}-1,
 \qquad
 A(2^n-1)=2^{n+1}-n-2.
\end{equation}
The top $Q$-degree of \cref{eq:EM} is
\begin{equation}\label{eq:B-M}
 B(M):=\sum_{m=1}^{M}ma_m
 =\frac12\sum_{j=0}^{J}2^jL_j(L_j+1).
\end{equation}
At a dyadic endpoint,
\begin{equation}\label{eq:B-dyadic}
 B(2^n-1)=4^n-(n+2)2^{n-1}.
\end{equation}
The classical limit of the normalizing product is
\begin{equation}\label{eq:multiscale-factorial}
 \lim_{Q\uparrow1}\frac{\Ecal_M(1;Q)}{(1-Q)^{A(M)}}
 =\prod_{m=1}^{M}m^{a_m}
 =\prod_{j=0}^{J}2^{jL_j}L_j!.
\end{equation}
Thus the finite determinant naturally defines a dyadic multiscale factorial.

\section{Mellin asymptotics of the $q$-Pochhammer tower}

Put $Q=\e^{-t}$, $t>0$, and
\[
 y=\log\frac1t,
 \qquad
 u_t=\log_2\frac1t=\frac y{\log2}.
\]
The identity \cref{eq:E-log} may be rewritten as
\begin{equation}\label{eq:E-Theta}
 -\log\Ecal(u;\e^{-t})
 =\sum_{\ell=1}^{\infty}\frac{u^\ell}{\ell}\Theta(\ell t).
\end{equation}
This transfers the heat complex dimensions directly to a tower of $q$-Pochhammer symbols.

\subsection{The subcritical regime $|u|<1$}

Let
\begin{equation}\label{eq:Pu}
 \Pcal_u(v)=\frac1{\log2}
 \sum_{k\ne0}\Gamma(\omega_k)\zeta(\omega_k)
 \Li_{1+\omega_k}(u)\e^{2\pi ikv}.
\end{equation}
For $u$ in a compact subset of the unit disk this series is absolutely convergent and analytic in both $u$ and $v$.

\begin{theorem}[Subcritical $q$-Euler asymptotic]\label{thm:E-subcritical}
Fix $\rho<1$.  Uniformly for $|u|\le\rho$, every $J\ge0$, and every $0<\varepsilon<1$,
\begin{align}\label{eq:E-subcritical-asymp}
 -\log\Ecal(u;\e^{-t})
 ={}&\frac{2\Li_2(u)}t
 +\frac{\log(1-u)}{2\log2}\,y
 +C(u)+\Pcal_u(u_t)\\
 &-\sum_{r=1}^{J}
 \frac{B_{2r}\Li_{2-2r}(u)}{(2r)!(2^{2r-1}-1)}t^{2r-1}
 +O_{J,\rho,\varepsilon}(t^{2J+1-\varepsilon}),\notag
\end{align}
where
\begin{equation}\label{eq:Cu}
 C(u)=-\frac1{2\log2}
 \left.\frac{\partial}{\partial s}\Li_s(u)\right|_{s=1}
 -C_\Theta\log(1-u).
\end{equation}
\end{theorem}

\begin{proof}
For $c>1$, Mellin inversion and absolute convergence give
\begin{equation}\label{eq:E-Mellin}
 -\log\Ecal(u;\e^{-t})
 =\frac1{2\pi i}\int_{c-i\infty}^{c+i\infty}
 \Gamma(s)\Li_{s+1}(u)\zeta_D(s)t^{-s}\dd s.
\end{equation}
At $s=1$ the residue is $2\Li_2(u)/t$.  At $s=0$, write
\[
 \zeta_D(s)=\frac{a_{-1}}s+a_0+O(s),
 \quad
 a_{-1}=-\frac1{2\log2},
 \quad
 a_0=-\frac14-\frac{\log(2\pi)}{2\log2}.
\]
The double-pole residue after multiplication by $t^{-s}$ is
\[
 \frac{\log(1-u)}{2\log2}y
 -\frac1{2\log2}\left.\partial_s\Li_s(u)\right|_{s=1}
 -C_\Theta\log(1-u).
\]
The poles $\omega_k$ give \cref{eq:Pu}.  At $s=-(2r-1)$, the residue equals the displayed Bernoulli--polylogarithm coefficient.  At negative even integers the trivial zeros of $\zeta$ cancel the poles of $\Gamma$.
\end{proof}

The expansion is nonuniform at $u=1$: $\Li_1(u)=-\log(1-u)$ diverges, and a second singular factor meets the pole at $s=0$.  The critical limit must be calculated separately.

\subsection{The critical product $u=1$}

Define
\begin{equation}\label{eq:PE}
 \Pcal_E(v)=\frac1{\log2}
 \sum_{k\ne0}\Gamma(\omega_k)\zeta(1+\omega_k)
 \zeta(\omega_k)\e^{2\pi ikv}.
\end{equation}
Let $\gamma_1$ be the first Stieltjes constant in
\[
 \zeta(1+s)=\frac1s+\gamma-\gamma_1s+O(s^2).
\]
Put
\begin{equation}\label{eq:CE}
 C_E=
 \frac{\zeta''(0)+\gamma_1+\gamma^2/2-\pi^2/12}{2\log2}
 -\frac{\log(2\pi)}4-\frac{\log2}{24}.
\end{equation}
Numerically, $C_E=-2.4612722039168654768\ldots$.

\begin{theorem}[Critical dyadic $q$-Pochhammer tower]\label{thm:E-critical}
For every $N>0$,
\begin{align}\label{eq:E-critical-asymp}
 -\log\Ecal(1;\e^{-t})
 ={}&\frac{\pi^2}{3t}
 -\frac{y^2}{4\log2}
 -\left(\frac14+\frac{\log(2\pi)}{2\log2}\right)y\\
 &+C_E+\Pcal_E(u_t)+\frac{t}{24}+O_N(t^N)
 \qquad(t\downarrow0).\notag
\end{align}
Equivalently, this is the asymptotic of
\begin{equation}\label{eq:E1-tower}
 \Ecal(1;\e^{-t})
 =\prod_{j=0}^{\infty}
 (\e^{-2^jt};\e^{-2^jt})_\infty.
\end{equation}
\end{theorem}

\begin{proof}
At $u=1$, \cref{eq:E-Mellin} becomes
\[
 \frac1{2\pi i}\int
 \Gamma(s)\zeta(s+1)\frac{\zeta(s)}{1-2^{-s}}t^{-s}\dd s.
\]
The pole at $s=1$ contributes $2\zeta(2)/t=\pi^2/(3t)$.  At $s=0$ the three factors $\Gamma(s)$, $\zeta(s+1)$, and $(1-2^{-s})^{-1}$ produce a triple pole.  Write
\[
 \Gamma(s)\zeta(1+s)
 =\frac1{s^2}+H_0+O(s),
 \quad
 H_0=-\gamma_1-\frac{\gamma^2}{2}+\frac{\pi^2}{12},
\]
and
\[
 \zeta_D(s)=\frac{a_{-1}}s+a_0+a_1s+O(s^2),
 \quad
 a_1=\frac{\zeta''(0)}{2\log2}
 -\frac{\log(2\pi)}4-\frac{\log2}{24}.
\]
The coefficient of $1/s$ after multiplication by $\e^{sy}$ is
\[
 \frac{a_{-1}}2y^2+a_0y+a_1+a_{-1}H_0,
\]
which is exactly the quadratic-log part and $C_E$.  The poles $\omega_k$ give \cref{eq:PE}.  The pole at $s=-1$ contributes $t/24$.  Every pole of $\Gamma$ at an integer $s=-n$, $n\ge2$, is canceled by a trivial zero of either $\zeta(s)$ or $\zeta(s+1)$.  The contour can therefore be moved arbitrarily far left, giving $O_N(t^N)$.
\end{proof}

\begin{remark}
The quadratic logarithm is a critical phenomenon.  For $|u|<1$ the origin is a double Mellin pole and produces a term linear in $y$; at $u=1$, the pole of $\zeta(s+1)$ raises the order to three and produces $-y^2/(4\log2)$.
\end{remark}

\section{A symmetric $q$-integer deformation of the sinc product}

\subsection{Symmetric $q$-integers}

For $0<Q<1$, define
\begin{equation}\label{eq:q-int}
 [m]_Q=\frac{Q^{-m/2}-Q^{m/2}}{Q^{-1/2}-Q^{1/2}}
 =Q^{-(m-1)/2}\frac{1-Q^m}{1-Q}.
\end{equation}
If $Q=\e^{-t}$, then
\begin{equation}\label{eq:q-int-sinh}
 [m]_Q=\frac{\sinh(mt/2)}{\sinh(t/2)}.
\end{equation}
Thus $[m]_Q\ge m$, $[m]_Q\to m$ as $Q\uparrow1$, and $[m]_Q$ grows exponentially in $m$ for fixed $Q<1$.  The multiplicative identity
\begin{equation}\label{eq:q-int-mult}
 [ab]_Q=[a]_Q[b]_{Q^a}
\end{equation}
will encode exact $q$-self-similarity.

By functional calculus, let $[D]_Q$ be diagonal with eigenvalues $[2^jn]_Q$.

\begin{definition}[Symmetric dyadic $q$-sinc determinant]\label{def:PhiQ}
Define
\begin{equation}\label{eq:PhiQ}
 \boxed{\ \Phi_Q(z)=\det(I-z^2[D]_Q^{-2})
 =\prod_{m=1}^{\infty}
 \left(1-\frac{z^2}{[m]_Q^2}\right)^{a_m}.\ }
\end{equation}
Also define the one-scale factor
\begin{equation}\label{eq:SQ}
 S_Q(z)=\det(I-z^2[N]_Q^{-2})
 =\prod_{n=1}^{\infty}\left(1-\frac{z^2}{[n]_Q^2}\right).
\end{equation}
\end{definition}
Both are entire functions of order zero for fixed $Q<1$.

\subsection{Exact $q$-scaling}

\begin{theorem}[$q$-renormalization equation]\label{thm:q-scaling}
The functional calculus version of $D\simeq N\oplus2D$ is
\begin{equation}\label{eq:qD-decomp}
 [D]_Q\simeq[N]_Q\oplus[2]_Q[D]_{Q^2}.
\end{equation}
Consequently,
\begin{equation}\label{eq:PhiQ-scale}
 \boxed{\ \Phi_Q(z)=S_Q(z)
 \Phi_{Q^2}\!\left(\frac{z}{[2]_Q}\right).\ }
\end{equation}
Iterating,
\begin{equation}\label{eq:PhiQ-scale-iter}
 \boxed{\ \Phi_Q(z)=\prod_{j=0}^{\infty}
 S_{Q^{2^j}}\!\left(\frac{z}{[2^j]_Q}\right).\ }
\end{equation}
As $Q\uparrow1$, this becomes the classical identity
\[
 \Phi(z)=\prod_{j=0}^{\infty}\sinc(\pi z/2^j).
\]
\end{theorem}

\begin{proof}
On the $j\ge1$ subspace, the eigenvalues are $[2(2^{j-1}n)]_Q$.  By \cref{eq:q-int-mult},
\[
 [2(2^{j-1}n)]_Q=[2]_Q[2^{j-1}n]_{Q^2}.
\]
This proves \cref{eq:qD-decomp}.  Apply determinant multiplicativity to get \cref{eq:PhiQ-scale}; repeated use and
\[
 [2^j]_Q=\prod_{h=0}^{j-1}[2]_{Q^{2^h}}
\]
give \cref{eq:PhiQ-scale-iter}.
\end{proof}

\subsection{Exact $q$-Pochhammer factorization}

Set
\begin{equation}\label{eq:eta-lambda}
 \eta=\frac{z(1-Q)}{\sqrt Q},
 \qquad
 \lambda+\lambda^{-1}=2+\eta^2.
\end{equation}
Either root may be chosen; every formula below is symmetric under $\lambda\leftrightarrow\lambda^{-1}$.

\begin{theorem}[$q$-Pochhammer representation of $\Phi_Q$]\label{thm:PhiQ-pch}
For every $m\ge1$,
\begin{equation}\label{eq:single-qfactor}
 1-\frac{z^2}{[m]_Q^2}
 =\frac{(1-\lambda Q^m)(1-\lambda^{-1}Q^m)}{(1-Q^m)^2}.
\end{equation}
Therefore
\begin{equation}\label{eq:PhiQ-E}
 \boxed{\ \Phi_Q(z)
 =\frac{\Ecal(\lambda;Q)\Ecal(\lambda^{-1};Q)}{\Ecal(1;Q)^2}\ }.
\end{equation}
Equivalently,
\begin{equation}\label{eq:PhiQ-tower}
 \Phi_Q(z)=\prod_{j=0}^{\infty}
 \frac{(\lambda Q^{2^j};Q^{2^j})_\infty
       (\lambda^{-1}Q^{2^j};Q^{2^j})_\infty}
      {(Q^{2^j};Q^{2^j})_\infty^2}.
\end{equation}
Likewise,
\begin{equation}\label{eq:SQ-pch}
 S_Q(z)=
 \frac{(\lambda Q;Q)_\infty(\lambda^{-1}Q;Q)_\infty}
      {(Q;Q)_\infty^2}.
\end{equation}
\end{theorem}

\begin{proof}
From \cref{eq:q-int},
\[
 \frac{z^2}{[m]_Q^2}
 =\frac{\eta^2Q^m}{(1-Q^m)^2}.
\]
The numerator after subtraction is
\[
 (1-Q^m)^2-\eta^2Q^m
 =1-(2+\eta^2)Q^m+Q^{2m}
 =(1-\lambda Q^m)(1-\lambda^{-1}Q^m),
\]
which proves \cref{eq:single-qfactor}.  Raise to the multiplicity $a_m$, multiply over $m$, and use \cref{eq:E-mult} and \cref{eq:E-tower}.
\end{proof}

\begin{resultbox}{Direct bridge to the classical sinc product}
Combining \cref{thm:PhiQ-pch} with the convergence theorem below gives the explicit limit
\[
 \Phi(z)=\lim_{Q\uparrow1}
 \prod_{j=0}^{\infty}
 \frac{(\lambda_Q(z)Q^{2^j};Q^{2^j})_\infty
       (\lambda_Q(z)^{-1}Q^{2^j};Q^{2^j})_\infty}
      {(Q^{2^j};Q^{2^j})_\infty^2},
\]
where $\lambda_Q(z)+\lambda_Q(z)^{-1}=2+z^2(1-Q)^2/Q$.  This is a direct $q$-Pochhammer representation of the infinite sinc product through a singular but controlled $Q\uparrow1$ limit.
\end{resultbox}

\subsection{Convergence to the Rvachev Fourier transform}

\begin{theorem}[Trace-norm and locally uniform limit]\label{thm:PhiQ-limit}
As $Q\uparrow1$,
\begin{equation}\label{eq:trace-norm-limit}
 \bigl\|[D]_Q^{-2}-D^{-2}\bigr\|_1\longrightarrow0.
\end{equation}
Consequently,
\begin{equation}\label{eq:PhiQ-limit}
 \boxed{\ \Phi_Q(z)\longrightarrow\Phi(z)\ }
\end{equation}
locally uniformly on $\C$.
\end{theorem}

\begin{proof}
The eigenvalues satisfy $[m]_Q\ge m$ and converge pointwise to $m$.  Hence
\[
 \bigl\|[D]_Q^{-2}-D^{-2}\bigr\|_1
 =\sum_{m=1}^{\infty}a_m
 \left(\frac1{m^2}-\frac1{[m]_Q^2}\right)
 \longrightarrow0
\]
by dominated convergence, dominated by $a_m/m^2$.  Continuity of the Fredholm determinant with respect to trace norm, uniformly when $z$ ranges over compact sets, proves the result.
\end{proof}

\subsection{Exact $q$-zero arithmetic and a Thue--Morse sign law}

\begin{theorem}[$q$-zero spectrum]\label{thm:q-zeros}
For fixed $0<Q<1$, the nonzero zeros of $\Phi_Q$ are exactly
\begin{equation}\label{eq:q-zero-set}
 z=\pm[m]_Q,\qquad m\ge1,
\end{equation}
with multiplicity $a_m=1+\nu_2(m)$.  Hence
\begin{equation}\label{eq:q-sign}
 \boxed{\ \sgn\Phi_Q(x)=(-1)^{s_2(M)}
 \quad\text{for}\quad [M]_Q<x<[M+1]_Q.\ }
\end{equation}
If $Q=\e^{-t}$, define
\begin{equation}\label{eq:MQ}
 M_Q(R)=\left\lfloor
 \frac2t\arsinh\!\left(R\sinh\frac t2\right)
 \right\rfloor.
\end{equation}
Then the number of nonzero zeros in $|z|\le R$, counted with multiplicity, is
\begin{equation}\label{eq:q-zero-count}
 \boxed{\ \mathcal N_{\Phi_Q}(R)
 =4M_Q(R)-2s_2(M_Q(R)).\ }
\end{equation}
For fixed $Q<1$, $\mathcal N_{\Phi_Q}(R)=O(\log R)$ and $\Phi_Q$ has order zero.  The limit $Q\uparrow1$ is a singular transition to the linear zero growth and order one of $\Phi$.
\end{theorem}

\begin{proof}
The sequence $[m]_Q$ is strictly increasing.  The product \cref{eq:PhiQ} gives the zero set and multiplicities.  Crossing the first $M$ positive zeros changes sign by $(-1)^{A(M)}=(-1)^{s_2(M)}$.  Formula \cref{eq:q-int-sinh} shows that $[m]_Q\le R$ exactly when
\[
 m\le\frac2t\arsinh\!\left(R\sinh(t/2)\right),
\]
which proves the count.  Exponential growth of $[m]_Q$ gives logarithmic counting and order zero.
\end{proof}

\subsection{The $q$-spectral zeta and coefficient recurrence}

For $\Re s>0$, define
\begin{equation}\label{eq:zetaDQ}
 \zeta_{D,Q}(s)=\Tr[D]_Q^{-s}
 =\sum_{m=1}^{\infty}\frac{a_m}{[m]_Q^s}.
\end{equation}
Let
\begin{equation}\label{eq:A-lambert}
 \Acal(x)=\sum_{m=1}^{\infty}a_mx^m
 =\sum_{j=0}^{\infty}\frac{x^{2^j}}{1-x^{2^j}}.
\end{equation}

\begin{proposition}[$q$-spectral Lambert expansion]\label{prop:zetaDQ}
For $\Re s>0$,
\begin{equation}\label{eq:zetaDQ-Lambert}
 \boxed{\ \zeta_{D,Q}(s)
 =(1-Q)^sQ^{-s/2}
 \sum_{k=0}^{\infty}\frac{(s)_k}{k!}
 \Acal\!\left(Q^{k+s/2}\right),\ }
\end{equation}
where $(s)_k=s(s+1)\cdots(s+k-1)$.  For $\Re s>1$,
\begin{equation}\label{eq:zetaDQ-limit}
 \zeta_{D,Q}(s)\longrightarrow\zeta_D(s)
 \qquad(Q\uparrow1).
\end{equation}
If
\[
 \Phi_Q(z)=\sum_{n=0}^{\infty}c_n(Q)z^{2n},
\]
then
\begin{equation}\label{eq:cQ-rec}
 c_0(Q)=1,
 \qquad
 c_n(Q)=-\frac1n\sum_{r=1}^{n}
 \zeta_{D,Q}(2r)c_{n-r}(Q).
\end{equation}
\end{proposition}

\begin{proof}
From \cref{eq:q-int},
\[
 [m]_Q^{-s}=(1-Q)^sQ^{-s/2}
 Q^{ms/2}(1-Q^m)^{-s}.
\]
Expand $(1-Q^m)^{-s}=\sum_{k\ge0}(s)_kQ^{mk}/k!$ and sum over $m$.  Dominated convergence using $[m]_Q\ge m$ proves \cref{eq:zetaDQ-limit}.  Finally,
\[
 \log\Phi_Q(z)=-\sum_{r\ge1}\frac{\zeta_{D,Q}(2r)}r z^{2r}
 \qquad(|z|<1),
\]
and logarithmic differentiation gives \cref{eq:cQ-rec}.
\end{proof}

\section{A base-$b$ extension}

The operator mechanism is not restricted to $2$.  Let $b\ge2$ be an integer and define
\begin{equation}\label{eq:Db}
 D_b e_{j,n}=b^jn\,e_{j,n},
 \qquad
 \Phi_b(z)=\det(I-z^2D_b^{-2})
 =\prod_{j=0}^{\infty}\sinc\!\left(\frac{\pi z}{b^j}\right).
\end{equation}
Let $\nu_b(m)=\max\{j:b^j\mid m\}$ and $s_b(M)$ be the sum of the base-$b$ digits of $M$.

\begin{theorem}[Base-$b$ spectral package]\label{thm:base-b}
One has
\begin{align}
 D_b&\simeq N\oplus bD_b,\label{eq:Db-scale}\\
 \operatorname{mult}_{D_b}(m)&=1+\nu_b(m),\label{eq:mb}\\
 \zeta_{D_b}(s)&=\frac{\zeta(s)}{1-b^{-s}},\label{eq:zetab}\\
 A_b(M):=\sum_{m\le M}(1+\nu_b(m))
 &=\sum_{j\ge0}\left\lfloor\frac{M}{b^j}\right\rfloor
 =\frac{bM-s_b(M)}{b-1}.\label{eq:Ab}
\end{align}
For symmetric $q$-integers,
\begin{equation}\label{eq:PhibQ}
 \Phi_{b,Q}(z)=\prod_{m\ge1}
 \left(1-\frac{z^2}{[m]_Q^2}\right)^{1+\nu_b(m)}
\end{equation}
satisfies
\begin{equation}\label{eq:PhibQ-scale}
 \Phi_{b,Q}(z)=S_Q(z)
 \Phi_{b,Q^b}\!\left(z/[b]_Q\right).
\end{equation}
The associated $q$-Euler determinant is
\begin{equation}\label{eq:Eb}
 \Ecal_b(u;Q)=\prod_{j=0}^{\infty}
 (uQ^{b^j};Q^{b^j})_\infty.
\end{equation}
\end{theorem}

\begin{proof}
All statements follow by splitting off the $j=0$ summand and relabeling $j\mapsto j-1$.  The digit formula uses
\[
 \sum_{j\ge1}\left\lfloor\frac{M}{b^j}\right\rfloor
 =\frac{M-s_b(M)}{b-1}.
\]
The $q$-scaling follows from $[bn]_Q=[b]_Q[n]_{Q^b}$.
\end{proof}

The dyadic case is distinguished because the parity of \cref{eq:Ab} collapses to
\[
 (-1)^{A_2(M)}=(-1)^{s_2(M)},
\]
which is exactly Thue--Morse.  Thus the Thue--Morse sign is not an accidental decoration: it is the parity character singled out by the base-$2$ spectral tower.

\section{Computational verification and reproducible checks}

\subsection{What was checked}

The following identities were independently evaluated at high precision:
\begin{enumerate}[leftmargin=2em]
\item the leading zero-jet formula \cref{eq:first-derivative-zero} for integers of several distinct $2$-adic valuations;
\item the factorization \cref{eq:PhiQ-E} by comparing the direct zero product and the quotient of dyadic $q$-Euler determinants;
\item convergence $\Phi_Q(z)\to\Phi(z)$ at real points inside and outside the first zero interval;
\item the heat expansion \cref{eq:heat-asymp}, including the coefficient $-1/12$ of $t$ and the first Fourier coefficient of $\Pcal_\Theta$;
\item the critical expansion \cref{eq:E-critical-asymp}, whose residual after the displayed $t/24$ term was below $6\times10^{-12}$ in the tested range $0.01\le t\le0.2$ with a finite Fourier truncation.
\end{enumerate}

\subsection{$q\uparrow1$ sample values}

\begin{center}
\small
\begin{tabular}{@{}c|ccc@{}}
\toprule
$Q$ & $\Phi_Q(0.3)$ & $\Phi_Q(1.2)$ & $\Phi_Q(2.7)$\\
\midrule
0.80 & $0.8316534376$ & $-0.08646646$ & $-0.03017980$\\
0.90 & $0.8241751473$ & $-0.07432843$ & $-0.01331576$\\
0.95 & $0.8204894677$ & $-0.06907163$ & $-0.00902981$\\
0.98 & $0.8182876536$ & $-0.06613499$ & $-0.00721005$\\
0.99 & $0.8175541599$ & $-0.06518866$ & $-0.00669771$\\
\midrule
$Q=1$ limit & $0.8168198779$ & $-0.0642567960$ & $-0.0062253251$\\
\bottomrule
\end{tabular}
\end{center}

The convergence is not claimed to be monotone in general; the table merely illustrates the trace-norm theorem.

\subsection{A practical evaluation hierarchy}

The new identities suggest several complementary algorithms:
\begin{enumerate}[label=(\alph*),leftmargin=2em]
\item \textbf{Near the origin:} use the coefficient recurrence \cref{eq:c-rec} or its $q$-version \cref{eq:cQ-rec}.
\item \textbf{Near an integer zero:} use the normalized local factor \cref{eq:local-factor}; this removes catastrophic cancellation and exposes the exact multiplicity.
\item \textbf{On compact sets away from zeros:} use the finite canonical product with the certified tail estimate \cref{eq:det-tail}.
\item \textbf{For $Q<1$:} use the $q$-Pochhammer quotient \cref{eq:PhiQ-tower}, exploiting fast exponential convergence at large scales.
\item \textbf{For $Q$ close to $1$:} combine the exact quotient with the critical asymptotic \cref{eq:E-critical-asymp} to cancel large denominator terms analytically before numerical evaluation.
\end{enumerate}

\section{Connections back to the Fabius and Thue--Morse theories}

\subsection{Two independent appearances of $q$-binomial structure}

The existing Fabius corpus derives base-$1/2$ Gaussian-binomial coefficients from finite dyadic coefficient extraction and exact values of $F$ at dyadic rationals.  The present report derives Gaussian-binomial coefficients from a different object: finite spectral determinants of the zero divisor of $\widehat\Up$.  The common mechanism is multiscale dyadic refinement:
\[
 \text{dyadic value formulas}
 \longleftarrow \text{refinement blocks}
 \longrightarrow \text{spectral multiplicity blocks}.
\]
The two coefficient systems are not identical, but they live in the same category of products indexed by powers of $2$.  This opens a concrete comparison problem: identify transforms that carry the exact-value $q$-binomial arrays into the spectral arrays $C_{M,k}(Q)$ of \cref{eq:CMk}.

\subsection{Thue--Morse as a determinant sign}

Equation \cref{eq:spectral-parity} can be rewritten
\begin{equation}\label{eq:tm-determinant-sign}
 \tm_M=\sgn\det\!\left(I-x^2D^{-2}\right)
 \qquad(M<x<M+1).
\end{equation}
Thus a binary automatic sequence is recovered from signs of a single entire Fredholm determinant.  The $q$-deformed version \cref{eq:q-sign} moves the zero locations but preserves the automatic sign word exactly.

\subsection{Moments and zeros from one spectrum}

The same spectrum produces both:
\begin{itemize}[leftmargin=2em]
\item the zero multiplicities through its eigenvalue degeneracies;
\item the probability cumulants through the traces $\Tr D^{-2r}$.
\end{itemize}
This is stronger than a formal analogy.  It says the local probability law of the Rvachev density and the global zero divisor of its Fourier transform are two evaluations of the same spectral data.

\subsection{Heat smoothing versus Fourier decay}

The active Fourier drafts study $\Phi(x)$ at large real $x$, where lacunary products create pointwise exceptional behavior and metric limit laws.  The heat trace $\Theta(t)$ studies the same multiplicities under exponential smoothing.  It is much more rigid: all oscillations are captured by the explicit vertical lattice $2\pi i\Z/\log2$.  A promising next step is to seek a transform connecting the log-periodic heat term to the periodic amplitudes that occur in endpoint asymptotics of $F$.

\section{Lean formalization roadmap}

The operator-theoretic presentation is conceptually compact, but a successful Lean formalization should proceed from finite identities to analytic limits.

\subsection{Layer 1: arithmetic and finite products}

The most accessible statements are:
\begin{enumerate}[leftmargin=2em]
\item $a_m=1+\nu_2(m)$ and the recursions $a_{2m}=a_m+1$, $a_{2m+1}=1$;
\item $\sum_{m\le M}a_m=2M-s_2(M)$;
\item the finite factorization \cref{eq:finite-pch};
\item the Gaussian-binomial coefficient formula \cref{eq:CMk};
\item the finite form of the single-factor identity \cref{eq:single-qfactor};
\item the base-$b$ digit identity \cref{eq:Ab}.
\end{enumerate}
These require no infinite-dimensional analysis.

\subsection{Layer 2: locally finite zero factorization (delivered through the leading jet)}

This layer is now delivered.  The formal proof first establishes the total
shifted-sinc law, then the global denominator-cleared factorization at
\(2^kq\), and finally the canonical signed-integer form.  A local analytic
unit gives the exact complex order and the first nonzero derivative.  The
remaining local-jet work is the normalized all-orders coefficient API:
package \(R_m\), its logarithmic derivatives, and
\cref{eq:all-jets} for every positive \(\ell\).

\subsection{Layer 3: infinite products and traces}

One can next formalize:
\begin{enumerate}[leftmargin=2em]
\item absolute convergence of $\sum_{j,n}(2^jn)^{-2}$;
\item the regrouped canonical product;
\item the Dirichlet-series identity $\zeta_D(s)=\zeta(s)/(1-2^{-s})$ on $\Re s>1$;
\item trace-norm convergence of the $q$-deformation, represented as convergence of summable diagonal sequences.
\end{enumerate}
Only after these are stable is it efficient to introduce Mathlib's operator determinant infrastructure, if the available API is sufficiently mature.

\subsection{Layer 4: Mellin asymptotics}

The heat and $q$-Euler expansions require contour shifting, residues, zeta continuation, and vertical estimates.  They are the deepest formalization target.  A useful intermediate milestone is a theorem schema: given a meromorphic Mellin transform with a finite specified pole set and suitable vertical decay, derive an asymptotic expansion.  The dyadic results would then become applications.

\section{Frontier questions}

\begin{problem}[Arithmetic of half-integer amplitudes]
Classify or sharply estimate the numbers $\Phi((2k+1)/2)$ that govern every leading integer-zero jet through \cref{eq:Cm}.  Determine whether their signs and $2$-adic or transcendence properties admit formulae parallel to the exact dyadic values of $F$.
\end{problem}

\begin{problem}[Comparison of the two Gaussian-binomial webs]
Construct an explicit transform between the coefficient arrays in the exact dyadic-value formulas for $F$ and the spectral coefficients $C_{M,k}(Q)$ in \cref{eq:CMk}.  The shared multiscale factorization suggests a common incidence algebra or a common family of refinement matrices.
\end{problem}

\begin{problem}[Critical double scaling]
Analyze $\Phi_Q(z)$ when $Q=\e^{-t}\uparrow1$ and $z$ grows with $t^{-1}$.  The fixed-$Q$ function has order zero and logarithmic zero growth, while the limit has order one and linear zero growth.  Identify the transition kernels and scaling regimes.
\end{problem}

\begin{problem}[Heat periodicity and endpoint periodicity]
The heat trace has the explicit periodic wave \cref{eq:Ptheta}; endpoint asymptotics of the Fabius function also contain periodic dependence on a dyadic logarithmic phase.  Determine whether both are projections of a single meromorphic two-variable transform.
\end{problem}

\begin{problem}[Twisted spectral traces]
Let $J e_{j,n}=\tm_n e_{j,n}$.  For $\Re s>1$,
\begin{equation}\label{eq:twisted-zeta}
 \Tr(JD^{-s})
 =\frac{1}{1-2^{-s}}
 \sum_{n=1}^{\infty}\frac{\tm_n}{n^s}.
\end{equation}
Study its continuation, functional relations, and whether its singularities encode known Thue--Morse Dirichlet-series phenomena in a spectral form.
\end{problem}

\begin{problem}[Transfer operators]
Relate the diagonal spectral model $D$ to the transfer operators used in the Fourier-decay corpus.  The former captures the exact zero divisor and smoothed scale statistics; the latter captures real-axis orbit statistics.  A common operator category could unify zeros, decay, and probabilistic limit laws.
\end{problem}

\begin{conjecture}[Universality of the symmetric $q$-deformation]
Among deformations that (i) preserve the multiplicity word $1+\nu_2(m)$, (ii) obey an exact two-scale factorization, (iii) have order zero for fixed $Q<1$, and (iv) converge locally uniformly to $\Phi$, the symmetric $q$-integer model \cref{eq:PhiQ} is unique up to an analytic reparameterization of $Q$ and a scalar rescaling of $z$.
\end{conjecture}

\section{Conclusion}

The infinite sinc product of Rvachev's up-function admits a spectral organization that is invisible when the product is treated only factor by factor.  The diagonal operator $D$ turns dyadic refinement into the unitary decomposition $D\simeq N\oplus2D$.  Its eigenvalue degeneracies are the integer-zero multiplicities; its finite spectral parity is the Thue--Morse sequence; its zeta function is $\zeta(s)/(1-2^{-s})$; its complex dimensions generate a log-periodic heat wave; and its power traces recover the probability cumulants.

The same operator admits two complementary $q$-constructions.  The determinant $\Ecal(u;Q)=\det(I-uQ^D)$ is a dyadic tower of $q$-Pochhammer symbols with exact finite Gaussian-binomial coefficients and explicit Mellin asymptotics.  Replacing each spectral integer by its symmetric $q$-integer produces $\Phi_Q$, an order-zero entire function with the same multiplicity and Thue--Morse sign word, an exact $q$-renormalization equation, and a locally uniform limit to the Rvachev Fourier transform.  Formula \cref{eq:PhiQ-tower} is the central new bridge: it connects $q$-Pochhammer symbols directly to the zero divisor of an infinite sinc product, while \cref{eq:PhiQ-scale} shows that the connection is dictated by self-similarity rather than by a formal change of notation.

The resulting program is concrete.  Its finite arithmetic core, total local
integer-zero factorization, exact complex order, and leading jet are now
machine checked.  The higher normalized jets remain frontier work; the
Mellin theory supplies exact human-readable asymptotics, and its unresolved
comparison with the exact dyadic-value $q$-calculus offers a sharply defined
next frontier.

\appendix
\clearpage
\begingroup
\small
\setlength{\tabcolsep}{5pt}
\sloppy
\section{Recursive TeX corpus ledger}\label{app:ledger}

The following ledger records the recursive TeX corpus used for the audit.  Directory prefixes are shortened relative to
\path{Analysis/FabiusFunction/docs/}.  For typographic wrapping, spaces in two historical directory names are displayed as underscores.

\subsection{Living and canonical documents: 18}

\begin{longtable}{@{}r P{0.80\textwidth}@{}}
\toprule
No. & Relative document\\
\midrule
\endfirsthead
\toprule
No. & Relative document\\
\midrule
\endhead
1 & \path{FabiusFunction_Mathematical_Glossary/FabiusFunction_Mathematical_Glossary.tex}\\
2 & \path{Fabius_Function_and_Rvachev_Up/Fabius_Function_and_Rvachev_Up.tex}\\
3 & \path{Non_Elementarity_of_the_Fabius_Function/Non_Elementarity_of_the_Fabius_Function.tex}\\
4 & \path{fabius_lean_walkthrough/fabius_lean_walkthrough.tex}\\
5 & \path{non-formalized-research-frontiers/non-formalized-research-frontiers.tex}\\
6 & \path{non-formalized-research-frontiers/Thue_Morse_Formula_Atlas/Thue_Morse_Formula_Atlas.tex}\\
7 & \path{.../rvachev_up_fourier_decay/Rvachev_Up_Fourier_Decay-2/Rvachev_Up_Fourier_Decay.tex}\\
8 & \path{.../rvachev_up_fourier_decay/Rvachev_Up_Fourier_Decay-3/Rvachev_Up_Fourier_Decay.tex}\\
9 & \path{.../Rvachev_Up_Fourier_Decay_Comparative_Audit/Rvachev_Up_Fourier_Decay_Comparative_Audit.tex}\\
10 & \path{.../Rvachev_Up_Fourier_Decay_Gentle_Guide/Rvachev_Up_Fourier_Decay_Gentle_Guide.tex}\\
11 & \path{.../Rvachev_Up_Fourier_Decay_Second_Wave_Audit/Rvachev_Up_Fourier_Decay_Second_Wave_Audit.tex}\\
12 & \path{.../asymptotic-decay-rate-of-an-infinite-product-of-sinc-functions/asymptotic-decay-rate-of-an-infinite-product-of-sinc-functions.tex}\\
13 & \path{.../rvachev_up_fourier_decay-1/rvachev_up_fourier_decay.tex}\\
14 & \path{.../rvachev_up_fourier_decay-4/rvachev_up_fourier_decay.tex}\\
15 & \path{.../rvachev_up_fourier_decay-5/rvachev_up_fourier_decay.tex}\\
16 & \path{.../rvachev_up_fourier_decay-6/rvachev_up_fourier_decay.tex}\\
17 & \path{.../rvachev_up_fourier_decay-7/rvachev_up_fourier_decay.tex}\\
18 & \path{.../rvachev_up_fourier_decay-8/rvachev_up_fourier_decay.tex}\\
\bottomrule
\end{longtable}

\subsection{Frozen historical archive: 32}

\begin{longtable}{@{}r P{0.80\textwidth}@{}}
\toprule
No. & Archived document directory / TeX base name\\
\midrule
\endfirsthead
\toprule
No. & Archived document directory / TeX base name\\
\midrule
\endhead
1 & \path{early-notes/Fabius_Asymptotic/Fabius_Asymptotic.tex}\\
2 & \path{early-notes/K-fold_summation_over_the_signed_Thue-Morse_sequence/K-fold_summation_over_the_signed_Thue-Morse_sequence.tex}\\
3 & \path{standalone-studies/Non_Elementarity_of_the_Fabius_Function/...tex}\\
4 & \path{standalone-studies/Small_Argument_Asymptotics/Small_Argument_Asymptotics.tex}\\
5 & \path{standalone-studies/fabius-inverse-dyadic-closed-form/fabius-inverse-dyadic-closed-form.tex}\\
6 & \path{primary-exposition/Fabius_Function_and_Rvachev_Up/...tex}\\
7 & \path{primary-exposition/Fabius_Function_and_Rvachev_Up-2/...tex}\\
8 & \path{primary-exposition/Fabius_Function_and_Rvachev_Up-3/...tex}\\
9 & \path{primary-exposition/Fabius_and_Rvachev_Up-2/Fabius_and_Rvachev_Up.tex}\\
10 & \path{primary-exposition/Fabius_and_Rvachev_up/Fabius_and_Rvachev_up.tex}\\
11 & \path{primary-exposition/Fabius_function_and_up_function/Fabius_function_and_up_function.tex}\\
12 & \path{primary-supplements/Fabius_Function_Missing_Parts/...tex}\\
13 & \path{primary-supplements/Fabius_Function_Missing_Parts-2/...tex}\\
14 & \path{q-calculus/Fabius_q_Special_Functions/Fabius_q_Special_Functions.tex}\\
15 & \path{q-calculus/fabius-q-special-functions/fabius-q-special-functions.tex}\\
16 & \path{lean-walkthrough/Fabius_Function_Lean_Walkthrough/...tex}\\
17 & \path{lean-walkthrough/fabius_lean_walkthrough/...tex}\\
18 & \path{glossary/FabiusFunction_Mathematical_Glossary/...tex}\\
19 & \path{glossary/FabiusFunction_Mathematical_Glossary-2/...tex}\\
20 & \path{research-frontiers/Fabius_Dyadic_Asymptotic_Bridge/...tex}\\
21 & \path{research-frontiers/Fabius_Dyadic_Formulae_and_Alternative_Representations/...tex}\\
22 & \path{research-frontiers/Fabius_Dyadic_Formulae_to_Asymptotics/...tex}\\
23 & \path{research-frontiers/Fabius_Dyadic_q_Connections/...tex}\\
24 & \path{research-frontiers/Fabius_Thue_Morse_Convergence_Rate/...tex}\\
25 & \path{research-frontiers/Fabius_Thue_Morse_Convergence_Rates/...tex}\\
26 & \path{research-frontiers/Primary_Exposition_Gap_Register/...tex}\\
27 & \path{research-frontiers/Repeated_Integration_and_Rvachev_Up/...tex}\\
28 & \path{research-frontiers/Rvachev_Up_from_Repeated_Integration/...tex}\\
29 & \path{research-frontiers/non-formalized-research-frontiers/...tex}\\
30 & \path{research-frontiers/Thue_Morse_Formula_Atlas-1/Thue_Morse_Formula_Atlas.tex}\\
31 & \path{research-frontiers/Thue_Morse_Formula_Atlas-2/Consolidated_Formula_Atlas_for_the_Thue_Morse_Sequence.tex}\\
32 & \path{research-frontiers/Thue_Morse_Formula_Atlas-3/Thue_Morse_Formula_Atlas.tex}\\
\bottomrule
\end{longtable}

\subsection{Imported paper sources: 7}

\begin{longtable}{@{}r P{0.80\textwidth}@{}}
\toprule
No. & Relative TeX source\\
\midrule
\endfirsthead
\toprule
No. & Relative TeX source\\
\midrule
\endhead
1 & \path{papers/AIF/AIF_2017__67_2_637_0.tex}\\
2 & \path{papers/AIF/AIF_2017__67_2_637_0-corrected.tex}\\
3 & \path{papers/Lacunary_products/Lacunary_products.tex}\\
4 & \path{papers/09-Function/09-Function.tex}\\
5 & \path{papers/157-Arithmetic-v3/157-Arithmetic-v3.tex}\\
6 & \path{papers/document/document.tex}\\
7 & \path{papers/ubiq15/ubiq15-corrected.tex}\\
\bottomrule
\end{longtable}

\endgroup

\section{Residue calculations for the critical $q$-Euler product}\label{app:residues}

For reference, the local series used in \cref{thm:E-critical} are
\begin{align*}
 \Gamma(s)&=\frac1s-\gamma+
 \left(\frac{\gamma^2}{2}+\frac{\pi^2}{12}\right)s+O(s^2),\\
 \zeta(1+s)&=\frac1s+\gamma-\gamma_1s+O(s^2),\\
 \frac1{1-\e^{-Ls}}&=\frac1{Ls}+\frac12+\frac{Ls}{12}+O(s^3),\\
 \zeta(s)&=-\frac12-\frac{\log(2\pi)}2s+\frac{\zeta''(0)}2s^2+O(s^3).
\end{align*}
Hence
\begin{align*}
 \Gamma(s)\zeta(1+s)
 &=\frac1{s^2}+
 \left(-\gamma_1-\frac{\gamma^2}{2}+\frac{\pi^2}{12}\right)+O(s),\\
 \frac{\zeta(s)}{1-\e^{-Ls}}
 &=-\frac1{2L}\frac1s
 -\frac14-\frac{\log(2\pi)}{2L}\\
 &\quad+
 \left(\frac{\zeta''(0)}{2L}
 -\frac{\log(2\pi)}4-\frac L{24}\right)s+O(s^2).
\end{align*}
Multiplication by $t^{-s}=\e^{sy}$ then yields the $y^2$, $y$, and constant terms in \cref{eq:E-critical-asymp}.

\section{Selected source map}

\begin{thebibliography}{99}

\bibitem{Repository}
V.~Reshetnikov,
\emph{ProveIt: Analysis/FabiusFunction documentation corpus},
GitHub repository, living tree and frozen archive,
\url{https://github.com/VladimirReshetnikov/ProveIt/tree/main/Analysis/FabiusFunction/docs},
accessed 27 August 2026.

\bibitem{Primary}
V.~Reshetnikov,
\emph{The Fabius Function and Rvachev's Up-Function},
source in the ProveIt documentation corpus,
\url{https://github.com/VladimirReshetnikov/ProveIt/blob/main/Analysis/FabiusFunction/docs/Fabius_Function_and_Rvachev_Up/Fabius_Function_and_Rvachev_Up.tex}.

\bibitem{Frontiers}
V.~Reshetnikov,
\emph{Non-formalized research frontiers for the Fabius function},
consolidated source in the ProveIt corpus,
\url{https://github.com/VladimirReshetnikov/ProveIt/blob/main/Analysis/FabiusFunction/docs/non-formalized-research-frontiers/non-formalized-research-frontiers.tex}.

\bibitem{Atlas}
V.~Reshetnikov,
\emph{Thue--Morse Formula Atlas},
source in the ProveIt corpus,
\url{https://github.com/VladimirReshetnikov/ProveIt/blob/main/Analysis/FabiusFunction/docs/non-formalized-research-frontiers/Thue_Morse_Formula_Atlas/Thue_Morse_Formula_Atlas.tex}.

\bibitem{AriasDefinition}
J.~Arias de Reyna,
\emph{An infinitely differentiable function with compact support: definition and properties},
2017, arXiv:1702.05442.

\bibitem{AriasArithmetic}
J.~Arias de Reyna,
\emph{Arithmetic of the Fabius function},
2017, arXiv:1702.06487.

\bibitem{AHL}
C.~Aistleitner, M.~Hofer, and G.~Larcher,
\emph{On parametric Thue--Morse sequences and lacunary trigonometric products},
2015, arXiv:1502.06738.

\bibitem{Toth}
L.~T\'oth,
\emph{Linear combinations of Thue--Morse Dirichlet series},
2022, arXiv:2211.13570.

\bibitem{AlloucheShallit}
J.-P.~Allouche and J.~Shallit,
\emph{Automatic Sequences: Theory, Applications, Generalizations},
Cambridge University Press, 2003.

\end{thebibliography}

\end{document}
